% TOP_PROBLEM: {"id": 30002149, "title": "Non-Disappearance of Nonnegative Scalar-Curvature Limits", "category_id": 16, "category": {"id": 16, "name": "physics", "display_name": "Mathematical Physics", "description": "Problems at the intersection of mathematics and physics.", "slug": "physics", "order_index": 16, "created_at": "2026-07-31T15:26:25.671Z"}, "status": "open", "statusQualification": "The convergence topology was recovered as intrinsic flat convergence (PDF page 50). The source reports non-disappearance under additional nonnegative-Ricci-curvature or uniform-contractibility assumptions; these are partial cases, not replacements for the scalar-curvature target."}
% FIRST_POSTED: 2026-10-11
% ENTRY_KIND: statement_only
% SUBMISSION: {"kind": "statement_only", "category": "top_problems", "problem_id": "30002149", "proposal_type": "established", "source": "Mathematical Aspects of General Relativity (2012), Oberwolfach Reports, DOI 10.4171/owr/2012/37, pp. 2269--2333. Primary report accessed 2026-10-11; relevant PDF pages 50, 51. https://doi.org/10.4171/owr/2012/37; Sormani, scalar curvature and intrinsic-flat convergence (catalogue-listed context; not independently consulted in this preparation). https://arxiv.org/abs/1606.08949", "status_evidence": "Canonical source status open retained. Supplied literature triage dated 2026-08-21: The no-disappearance statement remains a conjecture; intrinsic-flat no-cancellation is proved only under extra hypotheses such as Ricci lower bounds or uniform contractibility. Primary-report access checked 2026-10-11; this historical report and inherited triage do not certify current resolution of every component.", "status_checked_on": "2026-10-11", "duplicate_check": "Compared canonical ID, title and normalized target with the complete local UnsolvedMath registry; normalized exact-text matches: none; title matches: none. Related or overlapping canonical records: none identified in this batch. Primary source context and the existing OWR statement audit used for scope and duplicate review; an exhaustive mathematical-equivalence/literature search was not performed.", "attribution": "Original proposer not named in the supplied canonical record; attribution remains to the identified Oberwolfach contribution/report, not its editors as invented proposers. Explanatory statement preparation and source/semantic audit assisted by OpenAI GPT-6 Codex; no human mathematical authorship or independent proof verification is claimed."}
% SOURCE_STATEMENT: "Do sequences of manifolds with nonnegative scalar curvature and no closed interior minimal surfaces other than their boundaries avoid disappearing in the limit unless their volumes tend to zero?"
% Preparation and final self-audit: CONTRIBUTING.md prompts C and D; source/audit artifacts retained in the local physics_statements_2026_10_11 batch.
% Imported from: physics_statements.tex; UnsolvedMath ID 30002149
% !TeX program = lualatex
\documentclass[11pt]{article}
\usepackage[margin=25mm]{geometry}
\usepackage{fontspec}
\setmainfont{Latin Modern Roman}
\usepackage{amsmath,amssymb,mathtools}
\usepackage{unicode-math}
\setmathfont{Latin Modern Math}
\usepackage{xurl,hyperref}
\hypersetup{colorlinks=true,urlcolor=blue}
\setcounter{secnumdepth}{0}
\setlength{\parindent}{0pt}
\setlength{\parskip}{5pt}
\setlength{\emergencystretch}{3em}
\newcommand{\sref}[2]{\hyperlink{source:#1:#2}{[S#2]}}
\newcommand{\cataloguescope}[1]{\paragraph{Recorded target.}#1}
\newcommand{\R}{\mathbb R}
\newcommand{\C}{\mathbb C}
\newcommand{\N}{\mathbb N}
\newcommand{\Z}{\mathbb Z}
\newcommand{\Q}{\mathbb Q}
\newcommand{\F}{\mathbb F}
\newcommand{\eps}{\varepsilon}
\begin{document}
% UnsolvedMath ID 30002149; source catalogue code OWR-12009-006
\setcounter{section}{30002148}
\section{Non-Disappearance of Nonnegative Scalar-Curvature Limits}\label{30002149}
{\small UnsolvedMath ID 30002149 · Mathematical Physics}
\subsection{Definitions and mathematical statement}

\paragraph{Shared notation.}$\mathbb N=\{1,2,\ldots\}$; the standard number fields have their usual meanings. All additional conventions are stated locally.


\paragraph{Definitions and conventions.}
Scalar curvature is the trace $R_g$ of the Ricci tensor of a Riemannian metric. A minimal hypersurface has mean curvature zero; a closed interior minimal surface is compact, boundaryless and lies in the interior. In a geometric convergence theory, disappearance means that a region with nonzero volume has no surviving part in the limiting geometric object. Volume is the Riemannian measure of the manifolds in the sequence. The report uses intrinsic flat convergence of oriented integral current spaces; disappearance of the whole sequence means convergence to the zero current space, $d_{\mathcal F}(M_j,0)\to0$. The cited conjecture asks whether this can happen with volumes bounded away from zero under the stated curvature and minimal-surface conditions.

\paragraph{Target retained from the canonical source catalogue.}
\paragraph{Catalogue wording retained for traceability.}
Do sequences of manifolds with nonnegative scalar curvature and no closed interior minimal surfaces other than their boundaries avoid disappearing in the limit unless their volumes tend to zero?

\paragraph{Source scope and qualification.}
The convergence topology was recovered as intrinsic flat convergence (PDF page 50). The source reports non-disappearance under additional nonnegative-Ricci-curvature or uniform-contractibility assumptions; these are partial cases, not replacements for the scalar-curvature target.
\subsection{Short English statement}
Can positive-volume regions disappear in limits of nonnegative-scalar-curvature manifolds without interior minimal surfaces?
\subsection{Sources}
\begin{itemize}
\item[\hypertarget{source:30002149:1}{S1}] ULAM.AI and the UnsolvedMath Contributors, UnsolvedMath, record 30002149 (OWR-12009-006). Catalogue curation under CC BY 4.0; definitions and English exposition added locally.
\url{https://huggingface.co/datasets/ulamai/UnsolvedMath}.
\item[\hypertarget{source:30002149:2}{S2}] Mathematical Aspects of General Relativity (2012), Oberwolfach Reports, DOI 10.4171/owr/2012/37, pp. 2269--2333. Primary report accessed 2026-10-11; relevant PDF pages 50, 51.
\url{https://doi.org/10.4171/owr/2012/37}.
\item[\hypertarget{source:30002149:3}{S3}] Sormani, scalar curvature and intrinsic-flat convergence (catalogue-listed context; not independently consulted in this preparation).
\url{https://arxiv.org/abs/1606.08949}.
\end{itemize}
\label{end:30002149}
\end{document}
